\section{Algebraic replacement and specialization}
\label{sec:algebra}

The connection construction supplies formal potentials over $\mathbb C$,
algebraic restriction maps, and identities in their completed local rings.
We first record exactly the data used in this section. We then replace the
potentials to sufficiently high order along the wall ideals, correct the
restriction graph formally, and put all resulting identities over one
finitely generated coefficient ring. Only these new data will be
specialized to positive characteristic.

\begin{definition}[The algebraic input]
\label{alg:input-data}
Fix finitely many labels $1\leq i\leq d$, integers $n_i\geq0$, and
$D=\sum_i n_i$. Write
\[
 a_i=(a_{i,1},a_{i,2},a_{i,3}),\qquad
 y_i\in\mathbb C^{n_i},\qquad
 C=(C_i^+,C_i^-)_{i=1}^d\in\mathbb C^{2D},\qquad
 z=(a,t),\quad t\in\mathbb C^D.
\]
An algebraic germ in specified coordinates means a regular germ at the
specified complex point of an \'etale neighborhood of their affine
coordinate space, expanded in those coordinates. Finitely many such
germs can be placed in a common neighborhood. The following are the input
data established by the preceding geometric and connection constructions.
\begin{enumerate}
\item For each $i$ there is a local \'etale complex ring $A_i$ in
coordinates $(a_i,y_i)$, with completion
$B_i=\mathbb C[[a_i,y_i]]$, maximal ideal $\mathfrak m_i$, and an ideal
$\mathfrak k_i\subset\mathfrak m_i$. The formal potential
$P_i^{\mathrm{o}}\in\mathfrak m_i^3B_i$ satisfies, for every wall
coordinate $x$,
\[
 \partial_xP_i^{\mathrm{o}}\in\mathfrak k_iB_i,\qquad
 \partial_xP_i^{\mathrm{o}}\in A_i+\mathfrak k_i^bB_i
 \quad(b\geq1).
\]
The second assertion means that a representative in the \emph{same}
ring $A_i$ exists at every indicated ideal order; it does not assert that
the full gradient is algebraic.
\item There are centered algebraic germs $Y_i^\pm(z)$ and an algebraic
ideal $\mathfrak r$ in a local \'etale chart in $z$. All ideals
in the following identities are extended to
$B_X=\mathbb C[[z]]$, whose maximal ideal is $\mathfrak m$. Put
\[
 \phi_i^+(z)=(a_i,Y_i^+(z)),\qquad
 \phi_i^-(z)=(0,Y_i^-(z)),\qquad
 L_i^\pm=\phi_i^{\pm *}(\mathfrak k_i)B_X.
\]
Then $L_i^\pm\subset\mathfrak m$ and
\begin{equation}
 (L_i^\pm)^3\subset\mathfrak r.
 \label{alg:cube-input}
\end{equation}
\item With $Y=(Y_i^+,Y_i^-)_{i=1}^d$, define
\[
 \Psi^{\mathrm{o}}(a,C)
 =\sum_{i=1}^d\bigl(P_i^{\mathrm{o}}(a_i,C_i^+)
                    -P_i^{\mathrm{o}}(0,C_i^-)\bigr),
 \qquad
 p_i^{\mathrm{o}}(y_i)=\partial_{a_{i,1}}P_i^{\mathrm{o}}(0,y_i).
\]
Differentiation with respect to $C$ is performed before substitution.
The following identities hold in $B_X$:
\begin{equation}
 \bigl((\partial_C\Psi^{\mathrm{o}})(a,Y)\bigr)=\mathfrak r,
 \qquad \Psi^{\mathrm{o}}(a,Y)\in\mathfrak r^2,
 \label{alg:value-gradient-input}
\end{equation}
and, for every ordered triple $(i,j,k)$, with repetitions allowed,
\begin{equation}
 a_{i,1}a_{j,1}a_{k,1}\in\mathfrak r,\qquad
 p_i^{\mathrm{o}}(Y_i^-)
 p_j^{\mathrm{o}}(Y_j^-)
 p_k^{\mathrm{o}}(Y_k^-)\in\mathfrak r.
 \label{alg:triple-input}
\end{equation}
Here the original triple memberships and
\eqref{alg:cube-input} are consequences of the end-deviation cube law on
the original restriction graph.
\item Let $\sigma$ exchange $C_i^+$ and $C_i^-$ for every $i$, and set
$J=\partial_tY(0,0)$, a $2D$-by-$D$ matrix. The two tangent images are
complementary:
\begin{equation}
 \Delta=\det[\,J\ \ \sigma J\,]\ne0.
 \label{alg:tangent-input}
\end{equation}
\end{enumerate}
\end{definition}

No assertion about connections over a field of positive characteristic
will be needed. In particular the algebraic input is a package of
consequences of the connection argument, not an assumption that its
Chern--Simons integrals have algebraic power-series coefficients.

\subsection{A primitive modulo an arbitrary ideal power}

We use the following form of detection by a formal arc.

\begin{lemma}[Detection on a formal arc]
\label{alg:arc}
Let $S$ be a reduced complete Noetherian local $\mathbb C$-algebra with
residue field $\mathbb C$. Every nonzero element $b$ of its maximal
ideal has nonzero image under some continuous local
$\mathbb C$-algebra homomorphism $S\longrightarrow\mathbb C[[u]]$.
Consequently, if $I\subset B=\mathbb C[[v_1,\ldots,v_n]]$ is radical
and $P\in B$ satisfies $\partial_{v_j}P\in I$ for every $j$, then
$P-P(0)\in I$.
\end{lemma}

\begin{proof}
Choose a minimal prime of $S$ not containing $b$ and pass to the
corresponding complete local domain $S_0$. Its dimension $r$ is positive.
Choose $r-1$ elements $x_2,\ldots,x_r$ such that
$(b,x_2,\ldots,x_r)$ is primary to the maximal ideal. Such a choice is
possible because $b\ne0$ in a domain and
$\dim S_0/(b)\leq r-1$. A minimal prime $\mathfrak q$ over
$(x_2,\ldots,x_r)$ has height at most $r-1$, so it is not maximal.
It cannot contain $b$, since otherwise it would contain the displayed
maximal-primary ideal. In $S_1=S_0/\mathfrak q$ the ideal $(b)$ is
maximal-primary. The principal ideal theorem and nonmaximality of
$\mathfrak q$ give $\dim S_1=1$.
These dimension and height steps use the Noetherian local dimension
theorems \cite[Tags 00KQ, 0BBZ, and 00KW]{Stacks}.

The continuous map $\mathbb C[[b]]\longrightarrow S_1$ is injective:
a nonzero one-variable series is a power of $b$ times a unit. It is
finite. Indeed $S_1/(b)$ is a finite-dimensional complex vector space;
lift a basis, and successively expand remainders in powers of $b$.
Completeness, together with equivalence of the $(b)$-adic and maximal
ideal topologies, expresses every element as a $\mathbb C[[b]]$-linear
combination of those finitely many lifts.
The normalization of the complete discrete valuation ring
$\mathbb C[[b]]$ in the finite fraction-field extension
$\operatorname{Frac}(S_1)$ is a finite complete discrete valuation ring.
Its residue field is finite over $\mathbb C$, hence equals $\mathbb C$.
These conclusions use the complete-DVR extension theorem, applicable
because the fraction-field extension is finite and separable, and
completeness of finite algebras
\cite[Tags 09E8 and 0325]{Stacks}.
Choosing a uniformizer identifies the normalization with
$\mathbb C[[u]]$ as a $\mathbb C$-algebra: subtract successive complex
residue coefficients and use completeness
\cite[Tag 0C0S(2)]{Stacks}. Since $S_1$ is integral over
$\mathbb C[[b]]$, it embeds in this normalization, locally and
continuously. The resulting arc detects $b$.

For the consequence, suppose the class of $P-P(0)$ in $B/I$ is nonzero.
An arc as above detects it. If $v_j$ maps to $v_j(u)\in(u)$, the formal
chain rule gives
\[
 \frac{d}{du}P(v(u))
 =\sum_j(\partial_{v_j}P)(v(u))\,v_j'(u)=0.
\]
In characteristic zero a one-variable series with zero derivative is
constant. Its constant is $P(0)$, contradicting detection.
\end{proof}

\begin{lemma}[Algebraic replacement of a formal primitive]
\label{alg:replacement}
Let $A$ be the local ring at a complex point of an \'etale
neighborhood of affine $n$-space, with coordinates $v$ centered at that
point, and let $B=\widehat A=\mathbb C[[v]]$. Let $K$ be any ideal of
$A$. Suppose $P\in B$ satisfies
\begin{equation}
 \partial_{v_j}P\in KB,\qquad
 \partial_{v_j}P\in A+K^bB
 \quad\text{for every $j$ and every $b\geq1$}.
 \label{alg:replacement-hyp}
\end{equation}
For every $m\geq1$ there is $P_m\in A$ such that
$P-P_m\in K^mB$. No reducedness, primaryness, or equidimensionality
condition is imposed on $K$.
\end{lemma}

\begin{proof}
We first descend a finite Taylor tuple of $P$ along the reduced support
of $K$ to the Taylor tuple of a single element of $A$. We then choose
the length of that tuple so that a function with zero tuple belongs to
$K^m$, including at embedded primary components.

If $K=A$ there is nothing to prove. Otherwise put $I=\sqrt K$.
The reduced ring $A/I$ is essentially of finite type over $\mathbb C$,
hence excellent. Its completion is reduced: excellence makes the
completion map regular, and reducedness ascends along a regular map.
Also completion of the quotient identifies this completion with $B/IB$.
These results apply to a Noetherian local ring here
\cite[Tags 07QW, 07QU, 07GH, 07QK, and 00MA]{Stacks}.
Thus $IB$ is radical, and Lemma~\ref{alg:arc} gives
\begin{equation}
 P\equiv P(0)\pmod{IB}.
 \label{alg:constant-support}
\end{equation}

The coordinate derivations extend uniquely to $A$, because the chart is
\'etale, and they commute. For a positive integer $q$, let
\[
 T_q=(A/I)[s_1,\ldots,s_n]/(s_1,\ldots,s_n)^q,\qquad
 \tau_q(f)=\sum_{|\alpha|<q}
       \frac{\partial^\alpha f}{\alpha!}s^\alpha\pmod I.
\]
The product rule says that $\tau_q$ is a ring homomorphism; on the
coordinate ring it is the shift $v\mapsto v+s$.
Give $T_q$ its $A$-module structure \emph{through $\tau_q$}. Its
filtration by $S=(s_1,\ldots,s_n)$ has successive quotients that are
finite sums of $A/I$ with their ordinary coefficient action. The shifted
module $T_q$ is therefore finite over $A$.

For clarity, its completion is the usual finite Taylor target with a
shifted scalar action. If $\mathfrak m_A$ is the maximal ideal of $A$,
put $J_q=\tau_q(\mathfrak m_A)T_q$ and
$H_q=\mathfrak m_AT_q$, where the latter uses the ordinary coefficient
map. Since $J_q+S=H_q+S$ and $S^q=0$, expansion of ideal powers gives
\begin{equation}
 J_q^{\,N+q-1}\subset H_q^N,\qquad
 H_q^{\,N+q-1}\subset J_q^N
 \quad(N\geq1).
 \label{alg:cofinal-topologies}
\end{equation}
The identity of the underlying ring thus identifies the two completions,
and yields
\begin{equation}
 T_q\otimes_{A,\tau_q}B
 \cong (B/IB)[s_1,\ldots,s_n]/(s_1,\ldots,s_n)^q.
 \label{alg:taylor-completion}
\end{equation}
The $B$-action on the right is still the Taylor-shifted action.
Under this identification the completed map $\widehat\tau_q$ is the
formal Taylor map on $B$. Indeed $A$ is dense in $B$ and each of the
finitely many derivatives in the formula is continuous, losing at most
$|\alpha|$ powers of the maximal ideal; \eqref{alg:cofinal-topologies}
identifies that coefficientwise limit with the shifted-module limit.

Every coefficient of $\widehat\tau_q(P)$ has a representative in $A/I$.
The constant coefficient is supplied by
\eqref{alg:constant-support}. For $1\leq|\alpha|<q$, choose $j$ with
$\alpha_j>0$, put $\beta=\alpha-e_j$, and use
\eqref{alg:replacement-hyp} to choose $g_j\in A$ with
$\partial_{v_j}P-g_j\in K^qB$. The product rule gives the explicit
loss estimate
\begin{equation}
 \partial^\beta(K^qB)\subset K^{q-|\beta|}B\subset IB.
 \label{alg:derivative-loss}
\end{equation}
Hence $\partial^\alpha P/\alpha!$ is represented modulo $IB$ by
$\partial^\beta g_j/\alpha!$. These finitely many coefficients define
one element $x\in T_q$ whose image in
\eqref{alg:taylor-completion} is $\widehat\tau_q(P)$.

It is essential to descend the image condition for this \emph{whole}
tuple. Set $L_q=\ker\tau_q$, $U_q=\operatorname{im}\tau_q$, and
$Q_q=T_q/U_q$, using the shifted module structure throughout.
There are exact sequences of finite $A$-modules
\begin{equation}
 \begin{split}
 0&\longrightarrow L_q\longrightarrow A\longrightarrow U_q
    \longrightarrow0,\\
 0&\longrightarrow U_q\longrightarrow T_q\longrightarrow Q_q
    \longrightarrow0.
 \end{split}
 \label{alg:descent-sequences}
\end{equation}
Completion of finite modules is exact and agrees with tensoring by $B$;
$A\to B$ is faithfully flat
\cite[Tags 00MA--00MC]{Stacks}. The class of $x$ in $Q_q$ becomes zero
after this tensor product, so it was zero already. Thus
$x=\tau_q(P_q)$ for a single element $P_q\in A$. Exactness of the first
sequence additionally gives
\begin{equation}
 \ker\widehat\tau_q=L_qB,\qquad P-P_q\in L_qB.
 \label{alg:kernel-descent}
\end{equation}

The whole Taylor tuple has now descended to one algebraic primitive.
It remains to ensure that equality of that tuple measures the required
order along every component of the original ideal, including embedded
components. We now choose $q$ so that $L_q\subset K^m$. Take a finite primary
decomposition
\[
 K^m=\bigcap_{\lambda=1}^h Q_\lambda,\qquad
 \mathfrak p_\lambda=\sqrt{Q_\lambda},
\]
including all embedded components. For each $\lambda$, choose
$N_\lambda$ with
$\mathfrak p_\lambda^{N_\lambda}\subset Q_\lambda$; this follows by
taking powers of finitely many generators of the radical. For
$f\in L_q$ we have
\[
 \partial^\alpha f\in I\subset\mathfrak p_\lambda
 \qquad(|\alpha|<q).
\]
The Zariski--Nagata differential-power theorem gives
\begin{equation}
 f\in\mathfrak p_\lambda^{(q)}
   =\mathfrak p_\lambda^qA_{\mathfrak p_\lambda}\cap A.
 \label{alg:differential-order}
\end{equation}
Its hypotheses are satisfied at \emph{every} $\mathfrak p_\lambda$:
$A$ is essentially smooth over the characteristic-zero field
$\mathbb C$, and $\kappa(\mathfrak p_\lambda)/\mathbb C$ is separable.
The coordinate operators $\partial^\alpha/\alpha!$ of order less than
$q$ form an $A$-basis of the differential operators of that order, by
the \'etale coordinate description. These are precisely the
operators tested above; see
\cite[Theorem 3.6 and Lemma 2.3]{DSGJ2020}.

Choose $q\geq\max_\lambda N_\lambda$. Equation
\eqref{alg:differential-order} then puts $f$ in
$Q_\lambda A_{\mathfrak p_\lambda}$ for every $\lambda$. A primary
ideal contracts from localization at its radical:
\[
 Q_\lambda A_{\mathfrak p_\lambda}\cap A=Q_\lambda,
\]
since $sf\in Q_\lambda$ with $s\notin\mathfrak p_\lambda$ forces
$f\in Q_\lambda$. Intersecting gives $L_q\subset K^m$.
Together with \eqref{alg:kernel-descent} this proves the lemma.
\end{proof}

\begin{remark}[Normal order and embedded components]
\label{alg:embedded-example}
The differential test in \eqref{alg:differential-order} measures order
in the normal directions at each prime $\mathfrak p$, including
nonclosed primes. In the regular local ring $A_{\mathfrak p}$ the normal
cotangent space
$\mathfrak pA_{\mathfrak p}/(\mathfrak pA_{\mathfrak p})^2$
injects into $\Omega_{A/\mathbb C}\otimes_A\kappa(\mathfrak p)$.
The conormal sequence is injective here because smoothness and
separability give its image dimension $\operatorname{ht}\mathfrak p$,
the dimension of that normal space
\cite[Tags 00TT, 00NO, and 00TU]{Stacks}. The coordinate derivations
therefore span the dual normal directions over $\kappa(\mathfrak p)$.
A nonzero initial form of degree $j$ in the symmetric graded ring of
$A_{\mathfrak p}$ is detected by $j$ such directions: in characteristic
zero a nonzero degree-$j$ polynomial has a nonzero $j$-fold polarized
derivative. Lift their coefficients to $A_{\mathfrak p}$ and expand
the resulting composition of derivations. It is an
$A_{\mathfrak p}$-linear combination of coordinate derivatives of
orders at most $j$. Thus the vanishing of all derivatives of orders
less than $q$ rules out an initial degree $j<q$. This is the local
normal-order content of \eqref{alg:differential-order}; the possibly
singular quotient $A/\mathfrak p$ at the original closed point causes
no difficulty.

The embedded components also affect the required Taylor order.
For $A=\mathbb C[x,y,z]_{(x,y,z)}$ and $K=(x^2,xy)$ one has
\[
 \sqrt K=(x),\qquad
 K^2=(x^4,x^3y,x^2y^2)=(x^2)\cap(x,y)^4.
\]
Here $L_q=(x^q)$. Taylor order $q=4$ suffices for $K^2$, whereas
$q=2$ or $3$ does not. The second component is primary for the embedded
prime $(x,y)$, whose residue field after localization is $\mathbb C(z)$.
The proof used neither a closed-point residue field at that prime nor
an identification of global symbolic and ordinary powers.
\end{remark}

\subsection{Preserving the equations and obtaining algebraic witnesses}

We apply the replacement lemma at ninth ideal order. Since each
pulled-back wall ideal has cube in $\mathfrak r$, the resulting value
error lies in $\mathfrak r^3$ and the derivative error lies in
$\mathfrak m\mathfrak r$. These two containments preserve the
square-value identity and, by Nakayama's lemma, the entire gradient
ideal. The following proposition also retains the two cubic laws and
chooses algebraic coefficients witnessing all these memberships.

\begin{proposition}
\label{alg:replacement-preserves}
For the data of Definition~\ref{alg:input-data}, one can replace each
$P_i^{\mathrm{o}}$ by an algebraic germ $P_i$ with
$P_i-P_i^{\mathrm{o}}\in\mathfrak k_i^9B_i$. The $P_i$ still have order
at least three. Define
\[
 \Psi(a,C)=\sum_i\bigl(P_i(a_i,C_i^+)-P_i(0,C_i^-)\bigr),
 \qquad p_i(y_i)=\partial_{a_{i,1}}P_i(0,y_i),
\]
and, as columns and scalars in $B_X$, set
\[
 v=(\partial_C\Psi)(a,Y),\qquad g=\Psi(a,Y),
 \qquad \nu_i=p_i(Y_i^-).
\]
Then $(v)=\mathfrak r$, $g\in\mathfrak r^2$, and both types of triple
products in \eqref{alg:triple-input} belong to $(v)$ with the new
$p_i$. Moreover, there are algebraic germs $b_{jl}$ and row vectors
$h_+^{ijk},h_-^{ijk}$, in one common local \'etale chart in $z$,
such that
\begin{equation}
 g=v^{\mathsf T}bv,\qquad
 a_{i,1}a_{j,1}a_{k,1}=h_+^{ijk}v,\qquad
 \nu_i\nu_j\nu_k=h_-^{ijk}v.
 \label{alg:algebraic-witnesses}
\end{equation}
Here $b$ is a $2D$-by-$2D$ matrix, with no symmetry requirement.
\end{proposition}

\begin{proof}
Apply Lemma~\ref{alg:replacement} with $K=\mathfrak k_i$ and $m=9$.
Since $\mathfrak k_i\subset\mathfrak m_i$, the replacement changes no
terms of degree less than nine, in particular preserving cubic order.
For $E_i=P_i-P_i^{\mathrm{o}}$, differentiate first in the wall
variables and then substitute the original maps $\phi_i^\pm$.
The product rule gives $\partial E_i\in\mathfrak k_i^8B_i$, so
\begin{equation}
 \begin{split}
 \phi_i^{\pm *}E_i&\in(L_i^\pm)^9\subset\mathfrak r^3,\\
 \phi_i^{\pm *}(\partial E_i)&\in(L_i^\pm)^8
 \subset\mathfrak r^2\mathfrak m^2\subset\mathfrak m\mathfrak r.
 \end{split}
 \label{alg:ninth-order-estimates}
\end{equation}
The middle containment uses $(L_i^\pm)^3\subset\mathfrak r$ and
$L_i^\pm\subset\mathfrak m$, not an estimate for a derivative of a
pulled-back function.
The value estimate preserves $g\in\mathfrak r^2$.
If $v^{\mathrm{o}}=(\partial_C\Psi^{\mathrm{o}})(a,Y)$, then
$v-v^{\mathrm{o}}\in(\mathfrak m\mathfrak r)^{2D}$.
Thus $(v)\subset\mathfrak r$ and
$\mathfrak r=(v)+\mathfrak m\mathfrak r$. Nakayama's lemma applied
to the finite $B_X$-module $\mathfrak r/(v)$ proves $(v)=\mathfrak r$.

The positive triple products have not changed. For each negative
factor,
$\nu_i-p_i^{\mathrm{o}}(Y_i^-)\in\mathfrak r$ by the same derivative
estimate. Telescoping a product of three factors therefore preserves
its membership in $\mathfrak r=(v)$, including repeated labels.

All the new germs $P_i$, the original maps $Y_i^\pm$, their indicated
derivatives, and their compositions are algebraic. Place the finitely
many germs in $z$ in a common local \'etale ring $E$: pull back
the wall presentations along the maps $(a_i,Y_i^+)$ and $(0,Y_i^-)$,
take finite fiber products, and localize at the specified point.
Then $\widehat E=B_X$. Faithful flatness contracts the ideals
$(v)^2B_X$ and $(v)B_X$ to $(v)^2E$ and $(v)E$. Consequently the
three membership assertions admit the coefficients in
\eqref{alg:algebraic-witnesses} in $E$. We use these algebraic
coefficients, rather than any previously chosen formal witnesses.
\end{proof}

\subsection{One coefficient ring and a formal correction}

The replacement has produced algebraic witnesses for all the needed
ideal memberships. We can therefore select finitely many defining
equations and constants before choosing a coefficient ring. The next
steps put their coefficients in one ring and correct the graph using
only operations in that ring; this is what makes specialization
possible.

\begin{lemma}[Coefficients of finitely many algebraic germs]
\label{alg:finite-ring}
Given finitely many algebraic germs in finitely many coordinate spaces
over $\mathbb C$, and finitely many specified nonzero complex constants,
there is a finitely generated $\mathbb Z$-subalgebra $R\subset\mathbb C$
containing every coefficient of those germs and the inverses of the
specified constants.
\end{lemma}

\begin{proof}
Use finite \'etale presentations, near the chosen points, of the
form
\[
 \mathcal F(x,w)=0,\qquad w(0)=c,\qquad
 J_0=\partial_w\mathcal F(0,c),\quad\det J_0\ne0.
\]
Regular-function denominators can be included as additional variables
by equations $r f(x,w)=1$. Adjoin to $\mathbb Z$ the finitely many
polynomial coefficients, the entries of $c$, the inverses of all the
initial Jacobian determinants and denominator values, and the
specified inverses. This defines a finitely generated subring $R$.

Solve the equations by total degree in $x$. At each positive degree
the unknown homogeneous coefficient vector is multiplied by the
\emph{same} matrix $J_0$; the other terms are expressions in previously
determined coefficients. The inverse
$J_0^{-1}=\operatorname{adj}(J_0)/\det J_0$ already has entries in $R$.
Induction therefore puts every coefficient in $R$. This procedure
introduces no division by the degree or by a growing family of
factorials. Apply it to all the finite presentations at once.
\end{proof}

We also need a remainder formula valid over a coefficient ring that
need not contain $\mathbb Q$.

\begin{lemma}[Remainders without numerical denominators]
\label{alg:remainder}
Let $R$ be a commutative ring and
$\Phi(a,C)\in(a,C)^3R[[a,C]]$, where $C$ has $N$ entries. There is an
$N$-by-$N$ matrix $H(a,C,\delta)$ over $R[[a,C,\delta]]$, with every
entry in $(a,C,\delta)$, such that
\begin{equation}
 \Phi(a,C+\delta)
 =\Phi(a,C)+(\partial_C\Phi)(a,C)^{\mathsf T}\delta
        +\delta^{\mathsf T}H(a,C,\delta)\delta.
 \label{alg:remainder-formula}
\end{equation}
There is also a matrix $G(a,C,\delta)$ with positive-order entries
satisfying
\begin{equation}
 (\partial_C\Phi)(a,C+\delta)-(\partial_C\Phi)(a,C)
       =G(a,C,\delta)\delta.
 \label{alg:gradient-difference}
\end{equation}
Both constructions use only addition, multiplication, and integer
binomial coefficients. The difference of any formal scalar function
at $C+\delta$ and $C$ similarly factors as a row times $\delta$.
\end{lemma}

\begin{proof}
Expand a monomial $c_{\alpha\beta}a^\alpha C^\beta$ by the binomial
formula. After subtracting its constant and linear terms in $\delta$,
the remaining terms are
\[
 c_{\alpha\beta}\binom\beta\gamma
       a^\alpha C^{\beta-\gamma}\delta^\gamma,
 \qquad |\gamma|\geq2.
\]
Choose two occurrences among the $\delta$ factors by a fixed rule,
allowing the same index twice, and assign the coefficient after removing
those factors to the corresponding entry of $H$.
The remaining total degree is
$|\alpha|+|\beta|-2\geq1$. Summing defines $H$ coefficientwise;
each coefficient receives only finitely many contributions.
The other assertions follow by applying the one-factor version of the
same construction. For a first partial of $\Phi$, of order at least
two, its one-factor remainder has order at least one, giving the claim
for $G$.
\end{proof}

\begin{theorem}[Formal data after algebraic replacement]
\label{alg:formal-data}
The input in Definition~\ref{alg:input-data} yields a finitely generated
$\mathbb Z$-subalgebra $R\subset\mathbb C$, separated potentials
\[
 \begin{gathered}
 P_i\in(a_i,y_i)^3R[[a_i,y_i]],\qquad
 p_i(y_i)=\partial_{a_{i,1}}P_i(0,y_i),\\
 \Psi(a,C)=\sum_i\bigl(P_i(a_i,C_i^+)-P_i(0,C_i^-)\bigr),
 \end{gathered}
\]
and a centered formal map
$\widetilde Y(a,t)\in(z)R[[z]]^{2D}$ with the same first-order part as
the original $Y$. With
$\widetilde v=(\partial_C\Psi)(a,\widetilde Y)$, there are row vectors
$\widetilde h_+^{ijk},\widetilde h_-^{ijk}$ over $R[[z]]$ such that,
for every ordered triple of labels,
\begin{equation}
 \begin{split}
  \Psi(a,\widetilde Y)&=0,\\
  a_{i,1}a_{j,1}a_{k,1}&=\widetilde h_+^{ijk}\widetilde v,\\
  p_i(\widetilde Y_i^-)
  p_j(\widetilde Y_j^-)
  p_k(\widetilde Y_k^-)&=\widetilde h_-^{ijk}\widetilde v.
 \end{split}
 \label{alg:seven}
\end{equation}
The determinant
$\det[\partial_t\widetilde Y(0,0)\ \sigma\partial_t\widetilde Y(0,0)]$
is a unit of $R$. In particular $\widetilde Y(0,t)$ and
$\sigma\widetilde Y(0,t)$ parametrize smooth formal submanifolds
of the $C$-space with complementary tangent spaces. On both
submanifolds $F(C)=\Psi(0,C)$ vanishes identically.
The same ring $R$ contains the coefficients of the potentials, the
corrected map, and all the displayed membership witnesses.
\end{theorem}

\begin{proof}
Start with Proposition~\ref{alg:replacement-preserves}. Apply
Lemma~\ref{alg:finite-ring} to the unspecialized wall germs $P_i$, all
entries of $Y$, and all algebraic witnesses $b,h_+^{ijk},h_-^{ijk}$ in
\eqref{alg:algebraic-witnesses}. Include $\Delta^{-1}$ among the
specified inverses. Partial differentiation multiplies coefficients
by integers and centered substitution uses finite sums at each degree,
so $\Psi,p_i,v,g,\nu_i$ also have coefficients in this one ring $R$.
All identities in \eqref{alg:algebraic-witnesses} hold in $R[[z]]$,
since $R\subset\mathbb C$. Because $\Psi$ has order at least three and
$Y$ is centered, $v\in(z)^2R[[z]]^{2D}$.

Apply Lemma~\ref{alg:remainder} to $\Psi$, with $N=2D$, and look for
\[
 \widetilde Y=Y+Bv,
\]
where $B$ is an $N$-by-$N$ formal matrix. For $\delta=Bv$,
\eqref{alg:remainder-formula} and $g=v^{\mathsf T}bv$ give
\[
 \Psi(a,Y+Bv)
 =v^{\mathsf T}\bigl(b+B+B^{\mathsf T}H(a,Y,Bv)B\bigr)v.
\]
It suffices to solve the matrix equation
\begin{equation}
 B=-b-B^{\mathsf T}H(a,Y,Bv)B.
 \label{alg:matrix-recursion}
\end{equation}

Here is the coefficient construction over $R$.
Set $\mathcal T(B)=-b-B^{\mathsf T}H(a,Y,Bv)B$.
The entries of $H(a,Y,Bv)$ lie in $(z)$ for every $B$, since $a,Y,Bv$
are centered and $H$ has positive order. If $B-B'$ has entries in
$(z)^n$, changing either outer matrix factor changes
$B^{\mathsf T}H(a,Y,Bv)B$ only in $(z)^{n+1}$.
Changing the inner argument $Bv$ changes it in $(z)^{n+2}$ because
$v\in(z)^2$. Thus
\begin{equation}
 B-B'\in(z)^n\quad\Longrightarrow\quad
 \mathcal T(B)-\mathcal T(B')\in(z)^{n+1}.
 \label{alg:contraction}
\end{equation}
The constant coefficient is forced to be $B(0)=-b(0)$.
For every higher degree the coefficient on the right of
\eqref{alg:matrix-recursion} depends only on previously determined
coefficients of $B$ and on the already known coefficients of $b,H,Y,v$.
This determines a unique solution in $\operatorname{Mat}_N(R[[z]])$.
Equivalently, iterate $\mathcal T$ and use
\eqref{alg:contraction}. The recursion uses only ring operations;
there are no new numerical inverses. It proves the first identity in
\eqref{alg:seven} exactly.

We next transport the two triple laws to this corrected graph.
Equation~\eqref{alg:gradient-difference} gives
\begin{equation}
 \widetilde v=(1+M)v,\qquad
 M=G(a,Y,Bv)B\in(z)\operatorname{Mat}_N(R[[z]]).
 \label{alg:gradient-transport}
\end{equation}
Hence $A_1=1+M$ is invertible over the same ring, with
$A_1^{-1}=\sum_{n\geq0}(-M)^n$. Let
$\widetilde\nu_i=p_i(\widetilde Y_i^-)$.
The one-factor difference formula of Lemma~\ref{alg:remainder}, followed
by projection of $Bv$ to its $i$th negative block, gives a row
$q_i\in R[[z]]^{1\times N}$ satisfying
\begin{equation}
 \widetilde\nu_i-\nu_i=q_i v.
 \label{alg:negative-factor-transport}
\end{equation}
For all ordered triples the exact telescoping formula reads
\[
 \widetilde\nu_i\widetilde\nu_j\widetilde\nu_k-\nu_i\nu_j\nu_k
 =\bigl(\widetilde\nu_j\widetilde\nu_kq_i
       +\nu_i\widetilde\nu_kq_j+\nu_i\nu_jq_k\bigr)v.
\]
Consequently the required witnesses can be chosen as
\begin{equation}
 \begin{split}
 \widetilde h_+^{ijk}&=h_+^{ijk}A_1^{-1},\\
 \widetilde h_-^{ijk}&=
 \bigl(h_-^{ijk}+\widetilde\nu_j\widetilde\nu_kq_i
       +\nu_i\widetilde\nu_kq_j+\nu_i\nu_jq_k\bigr)A_1^{-1}.
 \end{split}
 \label{alg:witness-transport}
\end{equation}
These identities also hold when labels coincide. All coefficients stay
in $R$: the difference factorizations use integer coefficients, and
centered substitution and the geometric inverse series involve only
finitely many terms at each degree. This proves both remaining
identities in \eqref{alg:seven}.

Finally $Bv\in(z)^2$, so the first-order part of $Y$ and its paired
tangent determinant are unchanged. This determinant is a unit of $R$
by its choice. The map
\[
 (t,u)\longmapsto\widetilde Y(0,t)+\sigma J u
\]
has invertible linear term $[J\ \sigma J]$, and therefore has a formal
inverse over $R$. To see coefficient control directly, solve for each
homogeneous degree using this same invertible constant matrix. This
identifies the first graph with the coordinate subspace $u=0$; the
swapped graph is treated identically. They have complementary tangent
spaces. The value identity at $a=0$ makes $F$ vanish on the first;
$F(\sigma C)=-F(C)$ makes it vanish on the second.
\end{proof}

\begin{remark}
\label{alg:formal-not-geometric}
The correction changes the graph of evaluation, while retaining the
separated formula for $\Psi$. It may mix all graph variables, and no
geometric realization of $\widetilde Y$ is required. The original
end-deviation cube law is used only before correction. Afterwards,
\eqref{alg:gradient-transport}--\eqref{alg:witness-transport} prove the
needed laws entirely within the common formal coefficient ring.
\end{remark}

\subsection{Specialization}

\begin{corollary}[Formal data in large positive characteristic]
\label{alg:specialization}
For every sufficiently large rational prime $\ell$, the ring $R$ in
Theorem~\ref{alg:formal-data} has a homomorphism to a finite field
$\Bbbk$ of characteristic $\ell$. Reducing every coefficient by
such a homomorphism gives separated potentials of order at least three,
the identities \eqref{alg:seven} with their witnesses, and the two
complementary zero graphs over $\Bbbk$. In particular $\ell$ can
be chosen larger than any prescribed finite bound.
\end{corollary}

\begin{proof}
Because $R\subset\mathbb C$, the finite type
$\mathbb Q$-algebra $R_{\mathbb Q}=R\otimes_{\mathbb Z}\mathbb Q$
is nonzero. Choose a maximal ideal. The weak Nullstellensatz, or
Zariski's lemma, says that its residue field $L$ is a finite extension
of $\mathbb Q$. We obtain a homomorphism $R\to L$.
The images of the finitely many ring generators, including all
designated inverses, lie in $\mathcal O_L[1/N]$ for some positive
integer $N$: multiplying each algebraic number by a suitable nonzero
integer makes it integral. Thus the homomorphism factors as
\[
 R\longrightarrow\mathcal O_L[1/N]\longrightarrow
       \mathcal O_L/\mathfrak l=\Bbbk
\]
for any prime ideal $\mathfrak l$ above a rational prime $\ell\nmid N$.
The residue field is finite. Every designated determinant stays a
unit, because its inverse is already in $R$.

Coefficientwise reduction commutes with multiplication, centered
substitution, and formal differentiation; each coefficient involves
only a finite sum and differentiation multiplies by an integer.
It therefore preserves every identity in \eqref{alg:seven}, including
the membership coefficients, and the unit tangent determinant.
It also preserves the order-at-least-three condition and the separated
formula for $\Psi$. The formal inverse construction in the theorem
reduces over the same ring, giving the claimed smooth zero graphs.
Only the algebraic replacements and their subsequent formal
recursions are reduced. The original complex Chern--Simons series,
and all integration used to obtain them, are absent from this
specialization step.
\end{proof}
